RelaySpec reuses a pretrained speculative drafter when its target is fine-tuned or replaced by a different model. We train linear maps that transform the new target's hidden states into inputs the existing drafter can use, while keeping the target and the original drafter weights frozen. The drafter uses these mapped features to propose tokens, and the new target verifies them. This lets us adapt the drafter to a changed target without retraining its prediction network. Our main implementation and experiments use DFlash~\citep{chen2026dflash}, so we describe its architecture and training procedures first, then apply the same construction to EAGLE-3.

\subsection{Adapting the drafter's inputs}
\label{sec:method-inputs}
We retain the drafter's $K$ target-layer inputs and place one map before each original fusion block~\citep{chen2026dflash}, with $K=5$ for DFlash. Map $W_i$ connects the new target's feature stream $i$ to the corresponding drafter input and is applied at every token position. A bias-free $W_i\in\mathbb R^{d_r\times d_t}$ converts $h_{i,t}\in\mathbb R^{d_t}$ to the width $d_r$ the drafter expects, and the frozen fusion block $F_i\in\mathbb R^{d_d\times d_r}$ projects it to the draft width $d_d$. The $K$ streams give the context
\begin{equation}
 \widehat c_t=N\!\left(\sum_{i=1}^{K}F_i W_i h_{i,t}\right),
 \label{eq:mapped-context}
\end{equation}
where $N$ is the original RMSNorm~\citep{zhang2019rmsnorm}. Only $W_1,\ldots,W_K$ are trained; the target, fusion projection, normalization, draft transformer, token embeddings and vocabulary head all remain frozen.

After training, we multiply each learned matrix $W_i$ by its corresponding fusion block $F_i$ and join the $K$ products into one matrix:
\begin{equation}
 G=[F_1W_1\;\cdots\;F_KW_K]\in\mathbb R^{d_d\times Kd_t},
 \qquad \widehat c_t=N(G[h_{1,t};\ldots;h_{K,t}]).
 \label{eq:folded-fusion}
\end{equation}
Block multiplication gives $G[h_{1,t};\ldots;h_{K,t}]=\sum_i F_iW_ih_{i,t}$, so one multiplication by $G$ replaces the separate map and fusion operations in exact arithmetic, with the original normalization applied afterward (Figure~\ref{fig:relayspec-overview}). We keep the maps linear because a linear map both changes the feature width and folds into the fusion projection, and a nonlinear alternative with more parameters does not improve on it (Section~\ref{sec:ablations}).

\subsection{Training for a LoRA-adapted target}
\label{sec:method-auf}
\textbf{Initialization.} LoRA fine-tuning~\citep{hu2022lora} leaves the feature dimensions unchanged, so we initialize each map to the identity. Then $F_iW_ih_{i,t}=F_ih_{i,t}$ for every input, and inserting the maps does not disrupt the existing drafter's behavior with the adapted target; this preserves the starting behavior with the active LoRA target, not the behavior before fine-tuning. Training learns changes $\Delta_i=W_i-I$ that help the frozen drafter predict the adapted target's tokens.

\textbf{Training data.} We generate responses with the adapted target and save their tokens and hidden states. At sampled response positions the drafter receives one known token followed by 15 masked positions, matching the pretrained Qwen DFlash block of 16~\citep{chen2026dflash}, so we adapt the existing drafter without changing its prediction procedure. Each block sees target features only from before its starting position, and attention connects positions within a block but not separate blocks.

\textbf{Objective.} Because verification stops at the first rejected proposal, we supervise the correct prefix and the first incorrect prediction, excluding later positions: the accept-until-fail (AUF) objective~\citep{yang2026specauf}. For block $b$ and position $j$, with $q_{bj}$ the predicted distribution, $y_{bj}$ the saved target token and $m_{bj}$ a mask excluding padding and non-response positions,
\begin{equation}
 \begin{gathered}
 u_{bj}=\mathbf 1\!\left[\arg\max_v q_{bj}(v)=y_{bj}\ \text{or}\ m_{bj}=0\right],
 \qquad s_{bj}=\operatorname{stopgrad}\bigl({\textstyle\prod_{k<j}u_{bk}}\bigr),\\
 \mathcal L_{\mathrm{AUF}}
 =-\frac{\sum_{b,j}m_{bj}s_{bj}\log q_{bj}(y_{bj})}
 {\sum_{b,j}m_{bj}s_{bj}}.
 \end{gathered}
 \label{eq:auf-method}
\end{equation}
The product includes only earlier predictions, so the first failure still receives supervision. We recompute this mask from current predictions and hold it fixed during differentiation, gradients pass through the frozen drafter to update the maps, and we normalize the loss within each microbatch before averaging across accumulation steps and workers.

\subsection{Training for transfer across models}
\label{sec:method-mse}
\textbf{Initialization.} A drafter trained with one target expects that target's hidden states, and replacing the target changes both their dimensions and meaning. We therefore train the maps to reconstruct features from the original target, initializing them with independent Xavier-uniform values~\citep{glorot2010initialization}.

\textbf{Training data.} We generate responses with the new target, then run both targets causally on the same saved prompt and response to collect paired hidden states; the original target does not generate a separate response and is not used during adapted inference. With a shared tokenizer both receive identical token IDs, and when tokenizers differ we pair only positions corresponding to the same text prefix.

\textbf{Objective.} Let $h_{bti}$ and $h^{(o)}_{bti}$ be the paired new- and original-target features in stream $i$. The original features supply the context we reconstruct:
\begin{equation}
 c^{(o)}_{bt}=N\!\left(\sum_{i=1}^{K}F_i h^{(o)}_{bti}\right).
 \label{eq:original-context}
\end{equation}
For $B$ examples, let $S_b$ contain the selected valid prompt and response positions in example $b$. Using relative squared error, we minimize
\begin{gather}
 D(u,v)=\frac{\lVert u-v\rVert_2^2}{\lVert v\rVert_2^2+\varepsilon},
 \qquad \varepsilon=10^{-6},
 \label{eq:relative-error}\\
 \mathcal L_{\mathrm{MSE}}
 =\frac1B\sum_{b=1}^{B}\frac1{|S_b|}\sum_{t\in S_b}
 \left[\frac1K\sum_{i=1}^{K}D(W_i h_{bti},h^{(o)}_{bti})
 +D(\widehat c_{bt},c^{(o)}_{bt})\right].
 \label{eq:mse-method}
\end{gather}
The first term averages reconstruction errors across the $K$ feature streams and the second reconstructs the combined context received by the draft layers, with equal weight. We exclude padding and average positions within each example before averaging examples. Unlike AUF, this objective trains the maps without passing gradients through the draft transformer. Dividing by each original vector's squared norm makes the error relative to its magnitude. Reconstruction is a training objective, not a guarantee of accepted tokens, so we evaluate live decoding separately.

\subsection{Extension to EAGLE-3}
\label{sec:method-eagle}
The construction is not specific to DFlash. EAGLE-3 combines features from three target layers before autoregressive drafting~\citep{li2025eagle3}, so $K=3$. We place one map before each of those inputs, leave the original feature projection, prediction network and autoregressive drafting procedure unchanged, and fold the maps into that projection exactly as in Equation~\ref{eq:folded-fusion}. We pair each setting with the same objective we use for DFlash: identity initialization with token supervision for a LoRA-adapted target, applying the accept-until-fail rule along EAGLE-3's own recurrent path, and Xavier initialization with reconstruction for a replaced target. Appendix~\ref{app:eagle-method} gives the recurrent objective in full.

\subsection{Generation and verification}
\label{sec:method-generation}
The new target computes hidden states during prompt processing and proposal verification. We retain states from the committed prefix and map them into context for subsequent drafting. The frozen drafter proposes tokens, and the target verifies them using its existing decoding procedure. After verification, the decoder removes rejected speculative positions from the caches before the next drafting step. The original target is no longer needed to supply features during generation. Cross-tokenizer transfer additionally requires a proposal-conversion procedure, which the hidden-state maps do not provide.

Under greedy decoding, changing the drafter does not change the target's output provided block verification makes the same choices as sequential target evaluation on identical prefixes; Appendix~\ref{app:greedy} gives the argument and its assumptions.

